\chapter{ОБОБЩЕНИЕ И ОЦЕНКА РЕЗУЛЬТАТОВ}
\label{cha:results}

\section{Полнота решения задач этапа}

Все работы календарного плана первого этапа выполнены. Соответствие
запланированных работ и полученных результатов приведено в
таблице~\ref{tab:summary}.

\begin{table}[ht]
  \centering
  \caption{Соответствие работ календарного плана и результатов этапа}
  \label{tab:summary}
  \small
  \begin{tabular}{p{0.46\textwidth} p{0.22\textwidth} c}
    \toprule
    Работа календарного плана & Результат & Статус \\
    \midrule
    Загрузка и отрисовка моделей URDF & \texttt{modeuler} & выполнено \\
    Загрузка сцен SDF и объектов STL & \texttt{modeuler} & выполнено \\
    Генерация кинематической схемы из STEP & \texttt{kinechain} & выполнено \\
    Алгоритм симуляции без GPU & \texttt{dyphur} & выполнено \\
    Базовый графический интерфейс & \texttt{modeuler-win} & выполнено \\
    Предварительное тестирование и эталонные показатели & \texttt{dyphur} & выполнено \\
    Численное интегрирование ОДУ на GPU & \texttt{dyphur} & выполнено \\
    \bottomrule
  \end{tabular}
\end{table}

\section{Достоверность результатов}

Достоверность полученных результатов обеспечивается несколькими независимыми
механизмами:
\begin{itemize}
  \item функциональная корректность модулей подтверждена набором модульных
        тестов;
  \item физическая корректность ядра проверяется сопоставлением GPU-реализации с
        эталонной CPU-реализацией на идентичных сценариях;
  \item воспроизводимость подтверждена контролем побитового детерминизма
        (совпадение хешей состояния при повторных запусках);
  \item воспроизводимость среды обеспечена контейнеризацией сборки и фиксацией
        версий зависимостей.
\end{itemize}

\section{Сравнение с существующими решениями}

Выбранные алгоритмические решения соответствуют современному уровню: позиционная
динамика XPBD~\cite{macklin2016xpbd}, симплектическое интегрирование, иерархии
ограничивающих объёмов применяются в передовых системах симуляции. Ключевое
отличие разрабатываемого комплекса --- ориентация на \emph{переносимое}
GPU-ускорение (единый код для CUDA/HIP/CPU) при обеспечении детерминизма, а также
наличие средств автоматической подготовки моделей из конструкторской
документации (STEP $\rightarrow$ кинематическая схема), что в совокупности не
представлено в рассмотренных аналогах (раздел~\ref{cha:method}).

\section{Направления дальнейших работ}

По результатам этапа намечены следующие направления:
\begin{enumerate}
  \item масштабирование числа тел на порядки и дальнейшая оптимизация
        GPU-конвейера, в первую очередь решателя ограничений, занимающего
        основную долю времени кадра;
  \item интеграция подсистем (загрузка, кинематика, симуляция) в единый
        прикладной комплекс через графический интерфейс;
  \item расширение набора поддерживаемых типов суставов и геометрических
        примитивов при построении кинематики из STEP;
  \item верификация физической точности на задачах с аналитическим решением.
\end{enumerate}

Отрицательных результатов, требующих прекращения дальнейших исследований, не
получено. Полученные на этапе заделы достаточны для перехода к следующему этапу.
